% Chapter 6 — Physical Design

\chapter{Physical Design}
\label{ch:physical}

\section{What ``physical'' means}
Physical design is about \textbf{performance on real hardware}: which columns we index, which data types we pick, and how we refresh heavy summaries. Students often confuse this with ``logical'' design---remember: logical is ``what tables exist''; physical is ``how PostgreSQL finds rows fast.''

\section{Indexes we actually rely on}
Indexes are like book indexes: they speed up lookups but cost space and slightly slow down inserts. LearnTrack indexes match how the API filters data:

\begin{itemize}[leftmargin=*]
  \item \textbf{Courses by instructor.} Instructors list ``my courses'' with \texttt{courses.instructor\_id}. Index on \texttt{instructor\_id} avoids scanning every course.

  \item \textbf{Enrollments by course.} Roster and enrollment counts filter \texttt{enrollments.course\_id}.

  \item \textbf{Activity by user and time.} Analytics often asks ``what did this student do recently?'' Indexes on \texttt{activity\_log(user\_id)} and \texttt{activity\_log(event\_at DESC)} support that pattern.

  \item \textbf{Quiz attempts.} Reporting joins attempts by \texttt{user\_id} and \texttt{quiz\_id}; indexes there keep quiz history queries usable as data grows.

  \item \textbf{Partial index on published courses.} Listing the public catalogue only needs rows where \texttt{deleted\_at IS NULL} and often \texttt{is\_published}. A partial index shrinks the scanned set.
\end{itemize}

\section{Materialized views and their indexes}
Materialized views store \textbf{precomputed query results} on disk. We put unique indexes on natural keys such as (\texttt{user\_id}, \texttt{course\_id}) for performance summaries so:

\begin{itemize}[leftmargin=*]
  \item lookups by student+course are fast, and
  \item PostgreSQL can refresh some views \texttt{CONCURRENTLY} in setups that support it (less blocking during refresh).
\end{itemize}

\section{Data types (why they matter)}
\begin{itemize}[leftmargin=*]
  \item \textbf{TIMESTAMPTZ} stores ``when'' in absolute time; clients can display local zones without changing stored facts.

  \item \textbf{NUMERIC(5,2)} for percentages avoids floating-point rounding surprises in grade averages.

  \item \textbf{UUID} for \texttt{user\_id} aligns with Supabase authentication IDs.

  \item \textbf{INT / SERIAL} for course and content IDs keeps joins narrow and simple in demo queries.
\end{itemize}

\section{Trade-off summary}
We accept more disk space and occasional refresh time for materialized views because instructor dashboards would otherwise run expensive multi-join aggregates on every page load. Raw event tables stay authoritative; summaries can lag by minutes unless refreshed---a conscious engineering trade teachers should understand when demoing.
